% ============================================================
% Introduction
% ============================================================

\section{Introduction}

A well-calibrated forecaster states 80\% confidence intervals that contain
the true outcome 80\% of the time. Decades of research on human judgment
document systematic failure to achieve this standard: people's intervals are
too narrow, a phenomenon termed overconfidence in interval estimates
\citep{Alpert1982, Lichtenstein1982, Klayman1999, Soll2004}. The consequences
are severe. Overconfident interval estimates underlie miscalibrated risk
assessments in medicine \citep{ChristensenSzalanski1981}, finance
\citep{BenDavid2013}, engineering \citep{Flyvbjerg2002}, and strategic
planning \citep{Lovallo2003}.

Large language models (LLMs) are now routinely used as decision-support tools
in precisely these domains, yet whether their probabilistic forecasts are
calibrated remains an open question \citep{Kadavath2022, Lin2022, Xiong2024}.
Prior work has focused on factual question-answering confidence or binary
classification probabilities. Virtually no study has examined LLM calibration
on interval estimates for real-world quantitative forecasting, or asked how
miscalibration varies across prediction horizons and domains.

We address these questions with three preregistered studies. Study~1
establishes a baseline: are three frontier LLMs overconfident when predicting
next-day equity prices with stated confidence intervals? Study~2 extends to
nine prediction horizons (1--22 days) and asks: does miscalibration scale
with forecast horizon? Study~3 extends across six heterogeneous domains
(equities, commodities, cryptocurrency, foreign exchange, weather, NBA game
totals) and asks: is miscalibration domain-specific, and can its direction
reverse?

The human overconfidence literature consistently shows that miscalibration
is domain-specific: experts are most overconfident in their own fields of
ostensible competence \citep{Klayman1999, Soll2004}. It also shows that
overconfidence in interval estimates increases with task difficulty and
time horizon \citep{Griffin1992}. If LLMs inherit domain and horizon
structure from training data, they may replicate or amplify these asymmetries.

Our results confirm and extend the human literature. GPT-4 is the most
overconfident model in every study; Claude is the most calibrated. Study~2
reveals a model divergence at longer horizons: Claude and GPT-4 become
increasingly miscalibrated as the prediction window grows, while Gemini
maintains near-target coverage even at 22~days. Study~3 uncovers a
domain-confidence inversion---all models overconfident in structured financial
markets (equities, commodities) but underconfident in speculative markets
(cryptocurrency, forex)---a pattern absent from prior accounts of LLM
calibration. This paper contributes to the LLM calibration literature
\citep{Kadavath2022, Xiong2024} by focusing on interval estimates rather
than point probabilities and by using real-world outcomes rather than
knowledge-base queries.
